\chapter[Sermon 25. Here Is Declared What the Elders Did About the Throne, and How They Sang unto God a Song of Praise.]{Here Is Declared What the Elders Did About the Throne, and How They Sang unto God a Song of Praise.}
\label{sermon:025}
\markright{Sermon 25. Here Is Declared What the Elders Did About the Throne ...}

And when these beasts gave glory and honor, and thanks to the one who sits on the throne, who lives forever and ever, the twenty-four elders fell down before him who sat on the throne, and worshipped him who lives forever, and cast their crowns before the throne, saying, “You are worthy, Lord, to receive glory and honor and power, for you have created all things, and for your will’s sake they exist and were created.”

This most godly vision, well and rightly understood, and held in faithful memory, instructs us rightly in judging rightly the works of God, that we should fear God, be patient, and submit ourselves wholly to God, and give all glory unto him. For this is the very fruit that comes unto us, and the end of all things that here are spoken.

And by repeating what the beasts did, he indicates what the twenty-four elders did. By this, we are clearly taught what we owe to God, and how we should judge his works, and how we ought to conduct ourselves in this matter.

Those beasts, that is to say, the whole number of creatures, whose ministry God uses in the government of things, ascribe three things unto God sitting, that is to say, ruling and governing all things, to God I say living forever, that is to say, eternal, living, and giving or inspiring life into all things. First in deed glory, [illegible], which is a majesty, or great estimation, a reputation, worship, or good opinion: when we think well of God, protesting that there is nothing better than he, greater, more worthy, more just, more holy and more excellent. This glory are we always commanded to give him, and to esteem nothing in this world dearer and more precious than God. Secondly, they give to him honor [illegible]; in Greek, this signifies honor and price, and the due and bounden duty that we owe to any. We owe unto God reverence and submission, as to the supreme good, and the only and true lord of all. Paul in Romans 13, speaking of obedience due to the magistrate, says: “To whom you owe fear, give fear. And to whom you owe honor, give honor.” In the third place follows benediction, which he called [illegible], that is, thanksgiving, and praise. For we are commanded to praise all the works of the Lord, and to give thanks for the same. Job is said to have blessed or thanked God for the most grievous affliction that he sent him. For he said: “Like as it pleased the Lord, so hath it been done: the name of the Lord be blessed.” Whilst the beasts do attribute all these things to him that sits on the throne, by their example they teach us what we should do, namely, to give all these and singular things unto God. Which if we do, all murmuring shall cease, and disputations commenced of searching and examining the works of God through curiosity. With the laud and praise of the beasts is joined the hymn or song of the twenty-fourth. Elders. This is the church triumphant, the company of all Saints, Patriarchs, Prophets, Apostles, Martyrs, \&c., as I declared to you before. Mortal men have not here an example of some one saint, or wise man: but of all holy, godly, wise and worthy men. They have put off their flesh, and want affections and errors: they are therefore of uncorrupted judgements, so that there can be no more clear or pure examples ministered to us. Three or four things are taught us concerning these Elders, which they did or performed, not to every body, but to him that sits on the throne, and lives forever and ever. For so are the titles of God repeated, whereof is spoken before. We told you also that the seats of the Elders were set round about the Throne, in which they sat clothed with white raiment, crowned with crowns of gold, living with him that lives forever.

They first rise from their seats or chairs, and fall upon their knees or face before God. And in falling or kneeling, they show a submission and humility of mind, that we might learn with great humility and reverence to submit our souls and bodies to our God, submitting, I say, ourselves and all our things to his good will and pleasure. But if the blessed souls, now purified and already enjoying the sight of God, fall down before the Lord, what should a wretched, miserable, mortal, and sinful man do? Let him be ashamed of rebellion and slothfulness, which sees such great submission in the most noble and godly souls of heavenly dwellers.

Then the saints worship, and worship in deed none other, but him that sits on the throne, and lives forever, the Father, the Son, and the Holy Ghost, God three and one, everlasting and almighty. Therefore let us also worship this God, following the example of all saints. We worship God with external adoration, if we uncover our heads, kneel, and bow before him. In spirit and truth and with inward worship, if we depend wholly on him, consecrate ourselves wholly unto him, and wholly look upon him as one, the only, soul, incomprehensible, most wise, best, and greatest, most righteous and most merciful. And they that thus fall down before the throne of God, and so worship him, do not contend with God about his works; they do not expostulate with God impatiently, why he does this, and permits that?

\index{Apocalypse (Book of)}\index{Primasius}To all these things is added that they pluck the crowns from their heads and cast them away before the throne, in honor of him who sits there. This is not only a remarkable modesty but also an humble humility without parallel. Primasius, an interpreter of the Apocalypse, affirms that whatever virtue and dignity they possess is truly assigned to God. For to him is rightly attributed whatever is won or obtained; from him, the one who overcomes is aided. Thus he says. They also testify and signify that they would not assume any godly power, that they would not reign, that they would not, as counselors of God, give counsel to God or prescribe even the smallest thing to him, but to submit to God all power, all rule, and the entire government, themselves and all others to be governed. For they have experience and see no one in the world, neither in heaven nor on earth, who is wiser, mightier, greater, or who governs all things more faithfully, more diligently, more safely, and better. Let us therefore, O brothers, submit to the judgment of the saints, and let us consent with them in all things.

\index{Papacy}Yes, and with clear words they testify why they cast off their crowns. Not that they are ungrateful to God, and do not value his gifts highly; but because they openly acknowledge that all glory is due to him alone. Therefore, they are in complete agreement with the beasts and all of God’s creatures, and, singing a hymn to the highest ruler, they confess that he is worthy to receive glory. And he receives it, not as if he lacked it before, but so that it may seem unworthy if either they or any other creature were to claim as their own those things that belong to God alone. These things belong to no creature. And they praise God highly, whom they call their Lord and God. Some copies add, “who is holy.” For they agree in all things with the beasts, which also cried, “Holy, holy, holy, Lord God Almighty.” To him they also gave glory and honor, as mentioned before. So also the elders ascribe the same things to him now. Especially, they attribute power to God and take it from themselves. Why, then, do the Papists attribute power and activity to the saints in heaven? Which, nevertheless, here they plainly attribute to God alone. John and Peter, while they lived, did not take kindly to the people seeming to attribute some godly power to them. For when they had healed a man who had been lame before the temple, and the people were amazed, they said, “Men of Israel, why do you marvel at this? Or why do you look at us as if through our own power or holiness we had brought about this healing? The God of our fathers has done this.” But how much less should we now think that, being delivered from all corruption, they would desire any godly power to be given or divine honor attributed.

They also add or give a reason why they submit themselves, and all that they possess, to God, and attribute to him glory, honor, and power. For you have created all things, and by your will they are, and were created. This glory of God is wonderful and immeasurable. How great you are, and that all power and glory are due to you, appears from the making and creation of the universal world. No one was with you at the creation, no one gave you counsel about what or how you should do, no one helped you in the slightest. Who then should approach you to share in power? Who should glory before you, God and maker of all things? You alone made all things, alone preserve all, and alone govern all. You will, and they were made; you said, and they were created. It was enough to have said, it was enough to have willed. Indeed, all things today have their being through your will, without any effort or toil on your part. You govern all things in the best and most excellent order. This testifies to the wonderful course of the stars, the pleasant change of things, the most sweet and plentiful fruits. Who then would not gladly submit himself and all that he has to you and to your government, who would not commit all his possessions to you? Who would not acknowledge the power and glory to be yours? Let us mark these things with attentive minds, that we may also appear before God as we see the saints in heaven appear. God grant us this.
